\chapter{Dirichlet's Theorem}\label{ch:dirichlet}

\begin{center}\itshape
Every unbounded arithmetic progression whose first term and common difference are integers without common factor
contains infinitely many primes.\\
\upshape --- Title of Dirichlet's memoir (1837)
\end{center}

Count the primes up to $80$ before any theorem. There are $22$ of them. Nine are $1$ modulo $4$: $5,13,17,29,37,41,53,61,73$. Twelve are $3$ modulo $4$: $3,7,11,19,23,31,43,47,59,67,71,79$. The prime $2$ is in neither class. Figure~\ref{fig:primesmod4} is that list. Class $3$ leads by three.

\portrait{dirichlet}{Peter Gustav Lejeune Dirichlet (1805--1859). Anonymous portrait, before 1859. Public domain.}

In 1853 Chebyshev wrote to Fuss that the excess of primes $4n+3$ over primes $4n+1$ seemed to persist. Rubinstein and Sarnak quote the letter; this chapter does not treat the letter as a theorem. The race up to $60\,000$ is still lopsided (Figure~\ref{fig:race}). The weighted sums are not lopsided. This chapter proves that $\sum(\log p)/p$, split by residue class, gives every eligible class the same share. Chapter~\ref{ch:explicit} explains why the unweighted count can look biased for a long time.

Euclid's argument adapts to a few progressions: $4n+3$ (Volume~2, Exercise~2.1), $6n+5$, and, with more work, $kn+1$ for every $k$ (Volume~2, Exercise~5.7). For $5n+2$ no such polynomial trick is known. Dirichlet's method treats every progression at once. The ingredients are the $L$-functions of Chapter~\ref{ch:characters} and the fact $L(1,\chi)\ne0$. What remains is to sort the prime sums by characters. The result is quantitative. In the weighted sense of Mertens's first theorem, the primes are equally distributed among the $\varphi(q)$ classes prime to $q$.

\begin{think}
Up to $x=60\,000$ there are at least as many primes $\equiv3\pmod4$ as primes $\equiv1\pmod4$ at every $x$, except on the tiny interval $26861\le x<26863$ (Figure~\ref{fig:race}). Does that
contradict equal distribution? What would ``equally distributed'' have to mean for such a race to be compatible with it?
\end{think}

\begin{figure}[ht]
\centering
\fig{ch04-new-primes}
\caption{The $22$ primes up to $80$. A charcoal bar marks a prime $\equiv1\pmod4$, nine of them. A coral bar marks a prime $\equiv3\pmod4$, twelve of them. The prime $2$ has no bar. Class $3$ leads by three. Theorem~\ref{thm:dirichlet} is about a weighted sum, not about this lead.}
\label{fig:primesmod4}
\end{figure}

\section{The theorem}

Throughout, $q\ge1$ is fixed, $a$ is an integer prime to $q$, and the constants implied by $O(\cdot)$ may depend on $q$.

On 27 July 1837 Dirichlet read this theorem to the Berlin Academy. He was $32$. He told the Academy that the proof is not purely arithmetical. Part of it follows a continuous real variable, the number $s>1$ in $L(s,\chi)$. The argument below is that argument, with the characters of Chapter~\ref{ch:characters} in place.

A small weighted check, before the general sum. Up to $30$, the primes $\equiv1\pmod4$ are $5,13,17,29$, and
\[
\sum\frac{\log p}p=0.802
\]
to three decimals. The primes $\equiv3\pmod4$ are $3,7,11,19,23$, and the same sum is $1.153$. Half of $\log 30$ is $1.701$. Each class sits within $1$ of that main term. The theorem says the gap stays bounded for every $x$.

\begin{theorem}[Dirichlet, 1837]\label{thm:dirichlet}
For $x\ge2$,
\[
\sum_{\substack{p\le x\\p\equiv a\,(q)}}\frac{\log p}p=\frac1{\varphi(q)}\log x+O(1).
\]
In particular, there are infinitely many primes $p\equiv a\pmod q$.
\end{theorem}

Let $S(x,\chi)=\sum_{n\le x}\chi(n)\Lambda(n)/n$. The proof filters the class of $a$ by characters (Volume~3, Proposition~11.15):
\[
\sum_{\substack{n\le x\\n\equiv a\,(q)}}\frac{\Lambda(n)}n=\frac1{\varphi(q)}\sum_\chi\overline{\chi(a)}\,S(x,\chi).
\]
The principal character produces the main term. Every other character contributes $O(1)$.

\begin{lemma}\label{lem:principal}
$S(x,\chi_0)=\log x+O(1)$.
\end{lemma}
\begin{proof}
Theorem~\ref{thm:mertens1} says $\sum_{n\le x}\Lambda(n)/n=\log x+O(1)$. The sum $S(x,\chi_0)$ drops the terms with $n=p^k$ and $p\mid q$. Those terms add up to at most
\[
\sum_{p\mid q}\log p\sum_{k\ge1}p^{-k}=\sum_{p\mid q}\frac{\log p}{p-1},
\]
which depends only on $q$.
\end{proof}

\begin{lemma}\label{lem:chisums}
Let $\chi\ne\chi_0$. For $y\ge1$:
\[
\begin{gathered}\sum_{m\le y}\frac{\chi(m)}m=L(1,\chi)+O\Bigl(\frac1y\Bigr),\qq\sum_{m\le y}\frac{\chi(m)\log m}m=O(1).\end{gathered}
\]
\end{lemma}
\begin{proof}
Let $A(t)=\sum_{n\le t}\chi(n)$, so $|A|\le\varphi(q)$ by Proposition~\ref{prop:orth}. Take $f$ of class $C^1$, decreasing to $0$ on $[y,\infty[$. Lemma~\ref{lem:abelint} gives
\[
\sum_{y<n\le x}\chi(n)f(n)=A(x)f(x)-A(y)f(y)-\int_y^x A(t)f'(t)\,\dd t,
\]
and therefore $\bigl|\sum_{n>y}\chi(n)f(n)\bigr|\le 2\varphi(q)\,f(y)$. With $f(t)=1/t$, the tail of $L(1,\chi)$ after $y$ is at most $2\varphi(q)/y$. With $f(t)=(\log t)/t$, which decreases on $[3,\infty[$, the tail after any $y\ge3$ is at most $2\varphi(q)(\log y)/y$. The series $\sum\chi(m)(\log m)/m$ converges, so its partial sums stay bounded.
\end{proof}

\begin{proposition}\label{prop:nonprincipal}
Let $\chi\ne\chi_0$. Then $S(x,\chi)=O(1)$.
\end{proposition}
\begin{proof}
Since $\log=\Lambda*\mathbf1$ (Chapter~\ref{ch:mertens}) and $\chi$ is completely multiplicative,
\[
\chi(n)\frac{\log n}n=\sum_{dm=n}\frac{\chi(d)\Lambda(d)}d\cdot\frac{\chi(m)}m.
\]
Sum over $n\le x$ and group by $d$ (Lemma~\ref{lem:convsum}):
\[
\begin{aligned}\sum_{n\le x}\frac{\chi(n)\log n}n&=\sum_{d\le x}\frac{\chi(d)\Lambda(d)}d\sum_{m\le x/d}\frac{\chi(m)}m\eqbr=\sum_{d\le x}\frac{\chi(d)\Lambda(d)}d\bigl(L(1,\chi)+O(\tfrac dx)\bigr)\end{aligned}
\]
by Lemma~\ref{lem:chisums}. The error terms add up to $O\bigl(\psi(x)/x\bigr)$. Proposition~\ref{prop:chebyshev}(a) says this is $O(1)$. The left side is $O(1)$ by Lemma~\ref{lem:chisums}. Therefore $L(1,\chi)\,S(x,\chi)=O(1)$. Theorem~\ref{thm:lnonzero} says $L(1,\chi)\ne0$, so $S(x,\chi)=O(1)$.
\end{proof}

Figure~\ref{fig:twistedbars} is the sum $S(x,\chi_{-4})$ at eight values of $x$. The bars stay between $-0.200$ and $-0.304$. At $x=80$, one has $\log 80=4.382$. The principal sum is about the logarithm. The twisted sum is not.

\begin{figure}[ht]
\centering
\fig{ch04-new-bars}
\caption{The partial sums $S(x)=\sum_{n\le x}\chi_{-4}(n)\Lambda(n)/n$ at $x=10,20,\dots,80$, rounded to three decimals: $-0.200$, $-0.209$, $-0.206$, $-0.219$, $-0.258$, $-0.252$, $-0.248$, $-0.304$. Proposition~\ref{prop:nonprincipal} says this sum stays $O(1)$. At $x=80$ one has $\log 80=4.382$, so the principal sum is already much larger.}
\label{fig:twistedbars}
\end{figure}

\begin{proof}[Proof of Theorem~\ref{thm:dirichlet}]
Apply the filter, Lemma~\ref{lem:principal} and Proposition~\ref{prop:nonprincipal}. Since $\chi_0(a)=1$ and $|\chi(a)|=1$,
\[
\sum_{\substack{n\le x\\n\equiv a\,(q)}}\frac{\Lambda(n)}n=\frac1{\varphi(q)}\log x+O(1).
\]
The terms with $n=p^k$ and $k\ge2$ contribute at most $\sum_p\log p/(p(p-1))$, which converges, as in the proof of Theorem~\ref{thm:mertens1}. Removing them leaves the sum over primes and changes the total by $O(1)$. The right side tends to infinity, so the sum on the left has infinitely many terms.
\end{proof}

\begin{beauty}
Here is the line where $L(1,\chi)\ne0$ entered. The identity $\sum\chi(n)(\log n)/n=L(1,\chi)\,S(x,\chi)+O(1)$ says that the weighted primes, twisted by $\chi$, form a bounded sum whenever $L(1,\chi)\ne0$. If $L(1,\chi)$ vanished, the same identity with one extra logarithm would force $S(x,\chi)\sim-\log x$ (Exercise~\ref{ex:ifzero}). The primes would then avoid the classes where $\chi$ equals $1$ and crowd into the others. For the real character modulo $4$, one of the two odd classes would receive a bounded share of $\sum(\log p)/p$. Non-vanishing at $s=1$ and equal distribution of the weighted primes are two descriptions of one fact.
\end{beauty}

\section{Consequences}

Mertens's three theorems in Chapter~\ref{ch:mertens} now split by residue class. The extra ingredient is Theorem~\ref{thm:dirichlet}, and the summation by parts is the one already used for $\sum 1/p$.

\begin{corollary}\label{cor:recipprog}
There is a constant $M(q,a)$ such that, for $x\ge2$,
\[
\sum_{\substack{p\le x\\p\equiv a\,(q)}}\frac1p=\frac1{\varphi(q)}\log\log x+M(q,a)+O\Bigl(\frac1{\log x}\Bigr).
\]
\end{corollary}
\begin{proof}
Write $\sum_{p\le t,\,p\equiv a}(\log p)/p=(1/\varphi(q))\log t+R_a(t)$ with $R_a$ bounded. Repeat the proof of Theorem~\ref{thm:mertens2} with $f(t)=1/\log t$. The main term produces $(1/\varphi(q))(\log\log x-\log\log 2)$. The integral of $R_a(t)/(t\log^2 t)$ converges, and its tail is $O(1/\log x)$.
\end{proof}

Summing over the classes, $\sum_a M(q,a)=M-\sum_{p\mid q}1/p$. The classes that are not prime to $q$ contain only the primes that divide $q$ (Exercise~\ref{ex:constmod4}). Figure~\ref{fig:mertensprog} shows the two classes modulo $3$.

\begin{figure}[ht]
\centering
\fig{ch04-mertensprog}
\caption{$\sum_{p\le x,\,p\equiv a\,(3)}\frac{\log p}p-\frac12\log x$ for $a=1$ (charcoal) and $a=2$ (coral), $10\le x\le10^6$, in logarithmic scale. Both are bounded, as
Theorem~\ref{thm:dirichlet} says, and both converge, to $-1.188$ and $-0.511$ (at $x=10^6$ they have reached $-1.187$ and $-0.512$). The gap, $0.677$, is not a bias of large primes: $0.368$ of it is
$\frac{L'}L(1,\chi_{-3})$, which reflects the whole sequence of primes (the prime $2\equiv2$ among them), and $0.308$ is the difference between the higher-power contributions of the two classes, which a sum
over primes leaves out. Chapter~\ref{ch:pnt} shows that the limits exist.}
\label{fig:mertensprog}
\end{figure}

\section{Equal distribution, and in what sense}

Theorem~\ref{thm:dirichlet} gives each class prime to $q$ the share $1/\varphi(q)$ of the logarithmically weighted primes. A count of primes, one by one, needs more. Let $\pi(x;q,a)$ be the number of primes $p\le x$ with $p\equiv a\pmod q$. The prime number theorem for arithmetic progressions says
\[
\pi(x;q,a)\sim\frac1{\varphi(q)}\frac x{\log x}.
\]
That asymptotic follows from the method of Chapter~\ref{ch:pnt} applied to $L(s,\chi)$ on the line $\operatorname{Re}s=1$. Chapter~\ref{ch:pnt} indicates the changes. Even that asymptotic does not say which class leads at a given $x$. Figure~\ref{fig:race} shows a lead that lasts a long way.

\begin{figure}[ht]
\centering
\fig{ch04-race}
\caption{The race $\pi(x;4,3)-\pi(x;4,1)$ for $x\le60\,000$. The class $3$ is ahead or level everywhere except at one prime. There are ten ties, from $5$ to $26921$. At $x=26861$ the class $1$ leads by one, until the next prime $26863\equiv3\pmod4$. The remark below records a conditional logarithmic density of about $99.59\%$, due to Rubinstein and Sarnak. That percentage is not a theorem of this chapter. Chapter~\ref{ch:explicit} explains the bias.}
\label{fig:race}
\end{figure}

\begin{remark}
The squares of odd integers are all $\equiv1\pmod4$, and $\psi(x;4,1)$ counts squares of primes, while $\pi$ does not. Passing from the smooth weights of $\psi$ to the raw count $\pi$ removes a square contribution of size about $\sqrt x/\log x$ from the class $1$. Under the Riemann hypothesis for $L(s,\chi_{-4})$, the error in $\pi(x;4,a)$ is $O(\sqrt x\log x)$. That error is larger than the square contribution $\sqrt x/\log x$ by two powers of $\log x$. The lead is a bias of logarithmic density, not a pointwise comparison of those two sizes. Rubinstein and Sarnak (1994) computed the density, under further hypotheses. The figure $99.59\%$ is theirs. It is conditional, and it is not a theorem of this chapter. Chapter~\ref{ch:explicit} develops the bias.
\end{remark}

\section{What we proved, and what comes next}

For every $q$ and every $a$ prime to $q$, Theorem~\ref{thm:dirichlet} says
\[
\sum_{\substack{p\le x\\p\equiv a\,(q)}}\frac{\log p}p=\frac1{\varphi(q)}\log x+O(1).
\]
Equal distribution, in this chapter, means that weighted sum. Each eligible class gets the same share of $\sum(\log p)/p$, and each such class contains infinitely many primes. Corollary~\ref{cor:recipprog} splits $\sum 1/p$ in the same way, and Exercise~\ref{ex:progproduct} splits the product. The proof used orthogonality and $L(1,\chi)\ne0$ from Chapter~\ref{ch:characters}, and then repeated the summations of Chapter~\ref{ch:mertens}.

The weight $(\log p)/p$ does not say who is ahead at a given $x$. The same theorem does not give
\[
\pi(x;q,a)\sim\frac1{\varphi(q)}\frac x{\log x}.
\]
We do not have that count yet, even for $q=1$. The case $q=1$ is the prime number theorem.

Chapter~\ref{ch:pnt} proves it. The new step leaves the real axis and studies $\zeta$ on the line $\operatorname{Re}s=1$. One inequality between cosines shows that $\zeta$ has no zero on that line. The same argument, run with $L(s,\chi)$, then gives the asymptotic for every class.

\section{Exercises}

\begin{exercise}[$\star$]\label{ex:euclid}
Show that there are infinitely many primes $p\equiv5\pmod6$ and infinitely many primes $p\equiv1\pmod4$, by Euclid-style arguments.
(For the second, consider the prime factors of $(2p_1\cdots p_k)^2+1$.)
\end{exercise}

\begin{exercise}[$\star$]\label{ex:constmod4}
Show that $M(4,1)+M(4,3)=M-\frac12$, with $M$ the Meissel--Mertens constant, and more generally $\sum_{a}M(q,a)=M-\sum_{p\mid q}\frac1p$, the sum
running over the classes prime to $q$.
\end{exercise}

\begin{exercise}[$\star\star$]\label{ex:density}
\emph{Dirichlet density.} Show that $\sum_{p\equiv a\,(q)}p^{-\sigma}=\frac1{\varphi(q)}\log\frac1{\sigma-1}+O(1)$ as $\sigma\to1^+$, directly from Theorem~\ref{thm:lbasic} and
Theorem~\ref{thm:lnonzero}. This was Dirichlet's original formulation.
\end{exercise}

\begin{exercise}[$\star\star$]\label{ex:progproduct}
Show that there is a constant $C(q,a)>0$ with $\prod_{p\le x,\,p\equiv a\,(q)}\bigl(1-\frac1p\bigr)\sim C(q,a)(\log x)^{-1/\varphi(q)}$.
\end{exercise}

\begin{exercise}[$\star\star$]\label{ex:two}
\emph{Signature exercise.} Show that the primes $p$ for which $2$ is a quadratic residue modulo $p$ get half of the weight:
$\sum_{p\le x,\,(2/p)=1}\frac{\log p}p=\frac12\log x+O(1)$. Do the same for $-1$ and $-3$ in place of $2$.
\end{exercise}

\begin{exercise}[$\star\star$]\label{ex:ifzero}
Suppose, contrary to Theorem~\ref{thm:lnonzero}, that $L(1,\chi)=0$ for some $\chi\ne\chi_0$, and let $M_\chi(x)=\sum_{n\le x}\frac{\mu(n)\chi(n)}n$.
\begin{enumerate}[label=(\alph*)]
\item Show that $L'(1,\chi)M_\chi(x)=\log x+O(1)$.
\item Show that $S(x,\chi)=-L'(1,\chi)M_\chi(x)+O(1)$, and conclude that $S(x,\chi)=-\log x+O(1)$.
\end{enumerate}
\end{exercise}

\begin{exercise}[$\star$]\label{ex:count80}
Among the primes $p\le80$, how many satisfy $p\equiv1\pmod4$, and how many satisfy $p\equiv3\pmod4$? What is $\pi(80;4,3)-\pi(80;4,1)$? The prime $2$ belongs to neither count.
\end{exercise}

\begin{exercise}[$\star\star$]\label{ex:weight30}
Compute $\sum(\log p)/p$ over the primes $p\le30$ with $p\equiv1\pmod4$, and the same sum for $p\equiv3\pmod4$. Give each total to three decimal places, and compare both with $\frac12\log 30=1.701$.
\end{exercise}

\begin{exercise}[$\star\star$]\label{ex:mod8share}
The classes prime to $8$ are $1,3,5,7$. Deduce from Theorem~\ref{thm:dirichlet} that each of them satisfies
\[
\sum_{\substack{p\le x\\p\equiv a\,(8)}}\frac{\log p}p=\frac14\log x+O(1).
\]
Conclude that the primes with $\bigl(\frac2p\bigr)=1$, namely the classes $1$ and $7$ together (Volume~2, Corollary~6.4), receive half the weight, and that the single class $7$ already contains infinitely many primes.
\end{exercise}

\begin{exercise}[$\star\star\star$]\label{ex:schi4}
Let $\chi=\chi_{-4}$. Compute $S(30,\chi)=\sum_{n\le30}\chi(n)\Lambda(n)/n$ to three decimal places, by listing the prime powers. Using Proposition~\ref{prop:nonprincipal}, explain why this value may stay near $-0.2$. If instead $L(1,\chi)=0$, what size would Exercise~\ref{ex:ifzero} force on $S(30,\chi)$? Compare that size with $-\log 30=-3.401$.
\end{exercise}

\section{Hints}

\begin{hint}{ex:euclid}
For $5\pmod6$: consider $6p_1\cdots p_k-1$. For $1\pmod4$: if $p\mid m^2+1$ with $p$ odd, then $-1$ is a square modulo $p$.
\end{hint}
\begin{hint}{ex:constmod4}
Add the asymptotic formulas of Corollary~\ref{cor:recipprog} over $a$ and compare with Theorem~\ref{thm:mertens2}.
\end{hint}
\begin{hint}{ex:density}
Filter with characters. Use $\log L(\sigma,\chi)=\sum_p\chi(p)p^{-\sigma}+O(1)$, $\log L(\sigma,\chi_0)=\log\frac1{\sigma-1}+O(1)$, and that $L(\sigma,\chi)$ tends to a non-zero limit for $\chi\ne\chi_0$.
\end{hint}
\begin{hint}{ex:progproduct}
Take logarithms; $\log(1-\frac1p)=-\frac1p+O(\frac1{p^2})$.
\end{hint}
\begin{hint}{ex:two}
$\bigl(\frac2p\bigr)=\chi_8(p)$, the real character modulo $8$ with $\chi_8(n)=1$ for $n\equiv\pm1\pmod8$ (Volume~2, Corollary~6.4). The set of primes is the union of two
classes modulo $8$.
\end{hint}
\begin{hint}{ex:ifzero}
(a) Evaluate in two ways
\[
\sum_{n\le x}\frac{\mu(n)\chi(n)}n\sum_{m\le x/n}\frac{\chi(m)}m\log\frac x{nm}.
\]
(b) $\Lambda=\mu*\log$. In both, use Lemma~\ref{lem:chisums}, Exercise~\ref{ex:lderivative} and
\[
\sum_{m\le y}\frac{\chi(m)\log m}m=-L'(1,\chi)+O\Bigl(\frac{1+\log y}y\Bigr).
\]
\end{hint}

\begin{hint}{ex:count80}
List the primes. The residue modulo $4$ of an odd prime is $1$ or $3$.
\end{hint}
\begin{hint}{ex:weight30}
The primes $\equiv1$ are $5,13,17,29$. The primes $\equiv3$ are $3,7,11,19,23$. Round each $(\log p)/p$ to three decimals and add.
\end{hint}
\begin{hint}{ex:mod8share}
Apply Theorem~\ref{thm:dirichlet} with $q=8$, so $\varphi(8)=4$. Add the two classes $1$ and $7$.
\end{hint}
\begin{hint}{ex:schi4}
For $p\equiv3\pmod4$ the power $p^k$ contributes $(-1)^k(\log p)/p^k$. The powers $9$, $27$ and $25$ all lie at most $30$.
\end{hint}

\section{Solutions}

\begin{solution}{ex:euclid}
\idea A number of the shape $6P-1$ has a prime factor $\equiv5\pmod6$, because a product of primes $\equiv1\pmod6$ is $\equiv1$. For the class $1$ modulo $4$, the prime factors of $m^2+1$ are forced there by the quadratic character of $-1$.
\writeup Suppose $p_1,\dots,p_k$ were all the primes $\equiv5\pmod6$, and set $N=6p_1\cdots p_k-1$. No prime factor of $N$ is $2$ or $3$, so every prime factor is $\equiv\pm1\pmod6$. They are not all $\equiv1$, because $N\equiv5\pmod6$. None of them is a $p_i$, because $p_i$ divides $N+1$. That is a contradiction. For the class $1$ modulo $4$, set $N=(2p_1\cdots p_k)^2+1$. Then $N$ is odd and greater than $1$. A prime $p$ dividing $N$ is odd, is not one of the $p_i$, and $-1$ is a square modulo $p$. Euler's criterion $\bigl(\frac{-1}p\bigr)=(-1)^{(p-1)/2}$ forces $p\equiv1\pmod4$.

These arguments need a polynomial whose prime divisors lie in the target class, or in that class together with the class $1$. Such polynomials exist for the classes $a$ with $a^2\equiv1\pmod q$. That existence is Schur's theorem (1912). Murty (1988) proved the converse: no other class has such a polynomial. The progression $5n+2$ is out of reach of Euclid's method.
\end{solution}

\begin{solution}{ex:constmod4}
\idea The classes prime to $q$ partition the primes, once the finitely many primes that divide $q$ are removed.
\writeup Sum Corollary~\ref{cor:recipprog} over the $\varphi(q)$ classes prime to $q$. The result is
\[
\sum_{\substack{p\le x\\p\nmid q}}\frac1p=\log\log x+\sum_a M(q,a)+O\Bigl(\frac1{\log x}\Bigr).
\]
For $x\ge q$ the left side equals $\sum_{p\le x}1/p-\sum_{p\mid q}1/p$. Theorem~\ref{thm:mertens2} writes that as $\log\log x+M-\sum_{p\mid q}1/p+O(1/\log x)$. Compare constant terms. For $q=4$ the only prime divisor of $q$ is $2$, and $1/2$ is the term that is subtracted.
\end{solution}

\begin{solution}{ex:density}
\idea This is Theorem~\ref{thm:dirichlet} in the shape Dirichlet used. Weight the primes by $p^{-\sigma}$ and let $\sigma$ fall to $1$, instead of cutting the sum at $x$. The logarithm of $L(\sigma,\chi)$ is the twisted prime sum, up to a convergent error.
\writeup By Theorem~\ref{thm:lbasic}(a),
\[
\log L(\sigma,\chi)=\sum_p\chi(p)p^{-\sigma}+E(\sigma,\chi),
\]
where $\log$ means the series in that theorem and $|E(\sigma,\chi)|\le\sum_p\sum_{k\ge2}p^{-k}\le\sum_p 1/(p(p-1))<1$. Filter with characters:
\[
\sum_{p\equiv a}p^{-\sigma}=\frac1{\varphi(q)}\sum_\chi\overline{\chi(a)}\bigl(\log L(\sigma,\chi)-E(\sigma,\chi)\bigr).
\]
For $\chi\ne\chi_0$, one has $L(\sigma,\chi)\to L(1,\chi)\ne0$ by Theorem~\ref{thm:lnonzero}, and $\exp(\log L(\sigma,\chi))=L(\sigma,\chi)$. So $\log L(\sigma,\chi)$ is a continuous logarithm of a function with a non-zero limit, and it stays bounded as $\sigma\to1^+$. For $\chi_0$, the product $L(\sigma,\chi_0)=\zeta(\sigma)\prod_{p\mid q}(1-p^{-\sigma})$ is real and positive, and it is asymptotic to $(\varphi(q)/q)\cdot 1/(\sigma-1)$. The logarithm $\log L(\sigma,\chi_0)$ is the real logarithm, so it equals $\log 1/(\sigma-1)+O(1)$. The claim follows. Dirichlet stated the theorem in this form.
\end{solution}

\begin{solution}{ex:progproduct}
\idea As in Mertens's third theorem, the logarithm of the product is minus $\sum 1/p$, plus a convergent series of size $O(1/p^2)$. There is no analogue of $e^{-\gamma}$ that names the constant for a single class.
\writeup Let $h$ be the function from the proof of Theorem~\ref{thm:mertens3}, so $0\le h(p)=-\log(1-1/p)-1/p\le 1/p^2$. Then
\[
\log\prod_{\substack{p\le x\\p\equiv a}}\Bigl(1-\frac1p\Bigr)=-\sum_{\substack{p\le x\\p\equiv a}}\frac1p-\sum_{\substack{p\le x\\p\equiv a}}h(p).
\]
The second sum converges to some $H(q,a)$, with tail $O(1/x)$. Corollary~\ref{cor:recipprog} turns the logarithm into
\[
-\frac1{\varphi(q)}\log\log x-M(q,a)-H(q,a)+O\Bigl(\frac1{\log x}\Bigr).
\]
Exponentiate. The constant is $C(q,a)=e^{-M(q,a)-H(q,a)}$. Unlike $e^{-\gamma}$, these constants have no closed form in general. For $q=4$ they are computed numerically.
\end{solution}

\begin{solution}{ex:two}
\idea A real non-principal character detects quadratic residuacity, and Proposition~\ref{prop:nonprincipal} says that character contributes $O(1)$. The primes therefore split evenly, in the weighted sense, between residues and non-residues.
\writeup For odd $p$, one has $\bigl(\frac2p\bigr)=1$ if and only if $p\equiv\pm1\pmod8$, which is the condition $(1+\chi_8(p))/2=1$. Therefore
\[
\begin{aligned}\sum_{\substack{p\le x\\(2/p)=1}}\frac{\log p}p&=\frac12\sum_{2<p\le x}\frac{\log p}p+\frac12\sum_{p\le x}\frac{\chi_8(p)\log p}p\eqbr=\frac12\log x+O(1),\end{aligned}
\]
by Theorem~\ref{thm:mertens1} and Proposition~\ref{prop:nonprincipal}. The prime powers again contribute $O(1)$. For $-1$, one has $\bigl(\frac{-1}p\bigr)=\chi_{-4}(p)$ for odd $p$. For $-3$, quadratic reciprocity gives $\bigl(\frac3p\bigr)=\bigl(\frac p3\bigr)(-1)^{(p-1)/2}$, so
\[
\Bigl(\frac{-3}p\Bigr)=\Bigl(\frac{-1}p\Bigr)\Bigl(\frac3p\Bigr)=\Bigl(\frac p3\Bigr),
\]
because $\bigl(\frac{-1}p\bigr)=(-1)^{(p-1)/2}$ cancels the sign from reciprocity. For $p>3$ this is $\chi_{-3}(p)$. The same splitting applies. In general, if $b$ is not a square, the map $p\mapsto\bigl(\frac bp\bigr)$ is a non-principal real character modulo $4|b|$, evaluated at $p$, and half of the weighted primes have $b$ as a residue. The bookkeeping with reciprocity is in Ireland and Rosen, Chapter~5.
\end{solution}

\begin{solution}{ex:ifzero}
\idea If $L(1,\chi)=0$, the derivative $L'(1,\chi)$ takes the place that $L(1,\chi)$ held in Proposition~\ref{prop:nonprincipal}. The convolutions then produce a main term $\log x$, where the non-vanishing case had only a bounded sum.
\writeup Let $y\ge1$. By Lemma~\ref{lem:chisums} and its proof, $\sum_{m\le y}\chi(m)/m=O(1/y)$, because $L(1,\chi)=0$. Exercise~\ref{ex:lderivative} gives $L'(1,\chi)=-\sum_m\chi(m)(\log m)/m$, and the same Abel estimate yields
\[
\sum_{m\le y}\frac{\chi(m)\log m}m=-L'(1,\chi)+O\Bigl(\frac{1+\log y}y\Bigr).
\]
Hence
\[
\begin{aligned}\sum_{m\le y}\frac{\chi(m)}m\log\frac ym&=\log y\sum_{m\le y}\frac{\chi(m)}m\eqbr-\sum_{m\le y}\frac{\chi(m)\log m}m\eqbr=L'(1,\chi)+O\Bigl(\frac{1+\log y}y\Bigr).\end{aligned}
\]
(a) Let
\[
T=\sum_{n\le x}\frac{\mu(n)\chi(n)}n\sum_{m\le x/n}\frac{\chi(m)}m\log\frac x{nm}.
\]
Group by $k=nm$. The coefficient of $k^{-1}\log(x/k)$ is $\sum_{nm=k}\mu(n)\chi(n)\chi(m)=\chi(k)\sum_{n\mid k}\mu(n)$, which equals $1$ for $k=1$ and equals $0$ otherwise. So $T=\log x$. On the other hand, with $y=x/n$,
\[
T=L'(1,\chi)M_\chi(x)+O\Bigl(\sum_{n\le x}\frac1n\cdot\frac nx\Bigl(1+\log\frac xn\Bigr)\Bigr),
\]
and
\[
\frac1x\sum_{n\le x}\Bigl(1+\log\frac xn\Bigr)\le1+\frac1x\bigl(\lfloor x\rfloor\log x-\log\lfloor x\rfloor!\bigr)=O(1)
\]
by Proposition~\ref{prop:chebyshev}(c). Therefore $L'(1,\chi)M_\chi(x)=\log x+O(1)$.

(b) Since $\log=\Lambda*\mathbf1$, Möbius inversion gives $\Lambda=\mu*\log$. Thus $\chi(k)\Lambda(k)=\sum_{nm=k}\mu(n)\chi(n)\chi(m)\log m$. Lemma~\ref{lem:convsum} gives
\[
\begin{aligned}S(x,\chi)&=\sum_{n\le x}\frac{\mu(n)\chi(n)}n\sum_{m\le x/n}\frac{\chi(m)\log m}m\eqbr=-L'(1,\chi)M_\chi(x)+O(1),\end{aligned}
\]
with the same error estimate as in (a). Part (a) then yields $S(x,\chi)=-\log x+O(1)$. For $\chi_{-4}$ this would mean that $\sum_{p\le x,\,p\equiv1\,(4)}(\log p)/p$ stays bounded: the class $1$ would receive a bounded share of the weighted primes. Part (a) also forces $L'(1,\chi)\ne0$. A double zero at $s=1$ is impossible even before Theorem~\ref{thm:lnonzero} is used.
\end{solution}

\begin{solution}{ex:count80}
\idea An odd prime is $1$ or $3$ modulo $4$. Count the two lists in Figure~\ref{fig:primesmod4}.
\writeup The primes $p\le80$ with $p\equiv1\pmod4$ are $5,13,17,29,37,41,53,61,73$. There are $9$. The primes with $p\equiv3\pmod4$ are $3,7,11,19,23,31,43,47,59,67,71,79$. There are $12$. The prime $2$ is $\equiv2$. Therefore $\pi(80;4,3)-\pi(80;4,1)=12-9=3$.
\end{solution}

\begin{solution}{ex:weight30}
\idea Each term $(\log p)/p$ is small, and there are only nine odd primes up to $30$. Add them by class.
\writeup Rounded to three decimals, the class $1$ contributes
\[
\frac{\log5}5+\frac{\log13}{13}+\frac{\log17}{17}+\frac{\log29}{29}=0.322+0.197+0.167+0.116=0.802.
\]
The class $3$ contributes
\[
\frac{\log3}3+\frac{\log7}7+\frac{\log11}{11}+\frac{\log19}{19}+\frac{\log23}{23}=0.366+0.278+0.218+0.155+0.136=1.153.
\]
Half of $\log 30$ is $1.701$. The two sums differ from $1.701$ by $0.899$ and $0.548$. Both gaps are $O(1)$, already visible at $x=30$. Theorem~\ref{thm:dirichlet} says that for $q=4$ each gap stays bounded as $x$ grows.
\end{solution}

\begin{solution}{ex:mod8share}
\idea Four classes share the weighted primes. Adding two of them produces the half that Exercise~\ref{ex:two} found in one step. A single class still gets a quarter, and a quarter of $\log x$ tends to infinity.
\writeup Here $\varphi(8)=4$, and $1,3,5,7$ are the residues prime to $8$. Theorem~\ref{thm:dirichlet} gives each of them the sum $\frac14\log x+O(1)$. Add the classes $1$ and $7$:
\[
\sum_{\substack{p\le x\\p\equiv1\text{ or }7\,(8)}}\frac{\log p}p=\frac12\log x+O(1).
\]
Volume~2, Corollary~6.4 says $\bigl(\frac2p\bigr)=1$ for odd $p$ exactly in those two classes. That is the half-weight in Exercise~\ref{ex:two}, split one class further. The single class $7$ has weighted sum $\frac14\log x+O(1)$, which tends to infinity, so there are infinitely many primes $p\equiv7\pmod8$.
\end{solution}

\begin{solution}{ex:schi4}
\idea Only prime powers carry $\Lambda$. For $\chi_{-4}$ the prime $2$ drops out, primes $\equiv1\pmod4$ contribute with a plus, and primes $\equiv3\pmod4$ alternate with the exponent.
\writeup The contributions $n=p^k\le30$, rounded to three decimals, are $-0.366$ from $3$, $+0.122$ from $9$, $-0.041$ from $27$, $+0.322$ from $5$, $+0.064$ from $25$, $-0.278$ from $7$, $-0.218$ from $11$, $+0.197$ from $13$, $+0.167$ from $17$, $-0.155$ from $19$, $-0.136$ from $23$, and $+0.116$ from $29$. Their sum is $-0.206$. Proposition~\ref{prop:nonprincipal} says $S(x,\chi)=O(1)$ because $L(1,\chi_{-4})=\pi/4\ne0$. A value near $-0.2$ is what that bounded sum looks like at $x=30$. If $L(1,\chi)$ were $0$, Exercise~\ref{ex:ifzero} would give $S(x,\chi)=-\log x+O(1)$. At $x=30$ the main term is $-\log 30=-3.401$, about sixteen times the computed value $-0.206$.
\end{solution}

\section*{Sources and further reading}

\begin{itemize}
\item P.~G.~L.~Dirichlet, \emph{Beweis des Satzes, dass jede unbegrenzte arithmetische Progression \dots}, Abhandlungen der Königlichen
Preussischen Akademie der Wissenschaften (1837), 45--81. Read to the Academy on 27 July 1837.
\item H.~L.~Montgomery and R.~C.~Vaughan, \emph{Multiplicative Number Theory I}, Cambridge, 2007, Chapter~4 (primes in arithmetic
progressions, with Mertens-type estimates).
\item M.~R.~Murty, \emph{Primes in certain arithmetic progressions}, Journal of the Madras University (1988), 161--169 (which progressions admit
a Euclid-style proof).
\item M.~Rubinstein and P.~Sarnak, \emph{Chebyshev's bias}, Experimental Mathematics 3 (1994), 173--197. The conditional density $99.59\%$ is theirs, and so is the quotation of Chebyshev's letter of 1853 to Fuss.
\item A.~Granville and G.~Martin, \emph{Prime number races}, American Mathematical Monthly 113 (2006), 1--33. A survey of the races
between classes.
\item The first lead change at $26861$ is J.~Leech, \emph{Note on the distribution of prime numbers}, J.\ London Math.\ Soc.\ 32 (1957), 56--58.
\end{itemize}
